\chapter{Criterio de preparación para Tomo III}

\section{Paso al nivel metodológico}
\begin{definition}
Se dice que una pieza del canon intermedio está preparada para Tomo III cuando puede ser trasladada al nivel metodológico sin perder su herencia canónica, su función operativa ni la legibilidad de su recorrido interno.
\end{definition}

\begin{observation}
La preparación para Tomo III no convierte todavía a la pieza en demo auditada ni en claim final. Solo fija que ya alcanzó madurez suficiente para entrar en tratamiento metodológico y selectivo.
\end{observation}

\section{Condición mínima de preparación}
\begin{definition}
Una pieza está mínimamente preparada para Tomo III cuando concurren simultáneamente:
\begin{enumerate}
  \item dependencia canónica visible;
  \item nivel operativo suficientemente estabilizado;
  \item protocolización mínima legible;
  \item trazabilidad operativa suficiente;
  \item función reconocible dentro de un posible uso metodológico, demostrativo o selectivo.
\end{enumerate}
\end{definition}

\begin{observation}
No se exige todavía validación exhaustiva ni clausura definitiva de resultados. Basta con una preparación suficiente para no entrar al Tomo III como material opaco o prematuro.
\end{observation}

\section{Preparación y configuración}
\begin{definition}
Una configuración operativa está preparada para Tomo III cuando su función, su herencia y su modo de tratamiento pueden ser reconstruidos y reutilizados sin contradicción inmediata en un contexto metodológico.
\end{definition}

\begin{observation}
Esta cláusula evita que el Tomo III reciba configuraciones cuya identidad tratable todavía no esté estabilizada.
\end{observation}

\section{Preparación y familia}
\begin{definition}
Una familia operativa está preparada para Tomo III cuando sus miembros pueden ser presentados como conjunto disciplinado, con rasgos comunes legibles y diferencias internas no destructivas para el uso metodológico.
\end{definition}

\begin{observation}
La preparación de una familia no exige uniformidad absoluta. Exige solo cohesión suficiente para entrar en un marco auditado o selectivo.
\end{observation}

\section{Preparación y traducción}
\begin{definition}
Una traducción entre niveles está preparada para Tomo III cuando puede explicitar qué conserva, qué modifica y bajo qué condiciones sigue siendo metodológicamente utilizable en el nivel de llegada.
\end{definition}

\begin{observation}
Esta cláusula impide llevar al Tomo III traducciones cuya utilidad dependa de saltos implícitos o de pérdidas no declaradas.
\end{observation}

\section{Condición mínima de no preparación}
\begin{definition}
Se considera que una pieza no está preparada para Tomo III cuando falla de manera significativa en dependencia visible, estabilización operativa, protocolización mínima o trazabilidad suficiente.
\end{definition}

\begin{observation}
Esta condición protege al Tomo III frente a material todavía prematuro, híbrido o insuficientemente consolidado.
\end{observation}

\section{Cláusula de alcance}
\begin{observation}
El criterio de preparación para Tomo III fija una frontera mínima entre canon intermedio y tratamiento metodológico. No constituye todavía una teoría completa de auditoría, validación final o selección definitiva de claims.
\end{observation}

